% ============================================================================
% AOGS POSTER v2 (from-scratch redesign) — A0 portrait, 1189 x 841 mm.
% Design goals: NO dead space — slim header, tight class-level margins,
% merged blocks (5 + footer instead of 7), composite figures that fill
% their blocks edge-to-edge, references as a full-width footer strip.
% AOGS requirements kept: Abstract ID + title, presenter underlined+bold,
% availability window, contact info.
% Numbers: ALL from poster_numbers.tex (generated by
% code/pipeline/figures/make_poster_figures.py — never hand-typed).
% Build: /Library/TeX/texbin/pdflatex aogs_poster.tex
% ============================================================================
\documentclass[25pt, a0paper, portrait, margin=6mm, innermargin=10mm,
               blockverticalspace=2mm, colspace=10mm]{tikzposter}
\usepackage[T1]{fontenc}
\usepackage{newtxtext,newtxmath}
\usepackage[semibold]{sourcesanspro}   % Isotherm: quiet sans for the horizons
\usepackage{microtype}                 % crisp justification (protrusion + expansion)
\usepackage{graphicx}
\newcommand{\horAfif}{14.0}
\newcommand{\lossAfif}{1.16}
\newcommand{\svfAfif}{0.985}
\newcommand{\kdstarAfif}{4.60}
\newcommand{\horAsev}{10.1}
\newcommand{\lossAsev}{0.18}
\newcommand{\svfAsev}{0.991}
\newcommand{\kdstarAsev}{7.08}
\newcommand{\rmsehayAfif}{2.03}
\newcommand{\rmsecpAfif}{1.74}
\newcommand{\bestzcpAfif}{5/30}
\newcommand{\bestrhocpAfif}{1500/1700/1700}
\newcommand{\rmsecKAfif}{0.90}
\newcommand{\bestzcKAfif}{10/30}
\newcommand{\bestrhocKAfif}{1500/1700/1800}
\newcommand{\rmsehayAsev}{0.54}
\newcommand{\rmsecpAsev}{0.39}
\newcommand{\bestzcpAsev}{5/30}
\newcommand{\bestrhocpAsev}{1500/1700/1700}
\newcommand{\rmsecKAsev}{0.36}
\newcommand{\bestzcKAsev}{5/30}
\newcommand{\bestrhocKAsev}{1300/1500/1900}
\newcommand{\rmsehomAfif}{2.31}
\newcommand{\rmsehomAsev}{3.76}
\newcommand{\xferatAsev}{1.64}
\newcommand{\xferatAfif}{1.73}
\newcommand{\svfdtirAfif}{1.15}
\newcommand{\svfdtallAfif}{1.30}
\newcommand{\svfdtirAsev}{0.65}
\newcommand{\svfdtallAsev}{0.74}
\newcommand{\shadowcoolAfif}{2.22}
\newcommand{\shadowcoolAsev}{0.48}
\newcommand{\stabcKAfif}{13}
\newcommand{\bootcicKAfif}{0.26--1.08}
\newcommand{\stabcKAsev}{36}
\newcommand{\bootcicKAsev}{0.25--0.43}
\newcommand{\driftmax}{35}
\newcommand{\closurewin}{0.70}


% ── house palette (lunar.plotting.style) ────────────────────────────────────
\definecolor{coral}{HTML}{B85B3A}
\definecolor{teal}{HTML}{2A6478}
\definecolor{forest}{HTML}{3D6E4A}
\definecolor{char}{HTML}{2A2520}
\definecolor{cream}{HTML}{F7F5F1}
\definecolor{gridc}{HTML}{E8E5E0}

% Plain academic grammar (matches a standard conference-poster grid):
% white page, PLAIN BLACK section headers (no colored pill boxes), and
% teal reserved sparingly for the author bar, small highlight bands,
% and the Conclusions strip --- exactly the sparse-color convention of
% a typical AOGS/Imperial-style poster.
\definecolorstyle{lunarstyle}{}{
  \colorlet{backgroundcolor}{white}
  \colorlet{framecolor}{white}
  \colorlet{titlefgcolor}{white}
  \colorlet{titlebgcolor}{teal}
  \colorlet{blocktitlebgcolor}{white}
  \colorlet{blocktitlefgcolor}{char}
  \colorlet{blockbodybgcolor}{white}
  \colorlet{blockbodyfgcolor}{char}
  \colorlet{innerblocktitlebgcolor}{forest}
  \colorlet{innerblocktitlefgcolor}{white}
  \colorlet{innerblockbodybgcolor}{gridc}
  \colorlet{innerblockbodyfgcolor}{char}
  \colorlet{notefgcolor}{char}
  \colorlet{notebgcolor}{gridc}
}
\usetheme{Default}
\usecolorstyle{lunarstyle}
\useblockstyle{Basic}
\usetitlestyle{Empty}
\tikzposterlatexaffectionproofoff

% ── ISOTHERM design system ──────────────────────────────────────────────────
% Sections are horizons: a teal index chip + a quiet sans title struck by a
% single hairline. One horizon per stratum; hairlines divide, they never box.
% Blocks carry an EMPTY tikzposter title and place \horizon at the body top,
% so the header is fully under our control (weight, tracking, rule).
\newcommand{\isochip}[1]{%
  \colorbox{teal}{\makebox[7mm][c]{%
    \sffamily\bfseries\color{white}\rule[-1.4mm]{0pt}{5.6mm}#1}}}
\newcommand{\horizon}[2]{%
  \noindent\isochip{#1}\hspace{3.5mm}%
  {\sffamily\bfseries\fontsize{31}{35}\selectfont\color{char}#2}%
  \par\vspace{1.6mm}{\color{teal}\rule{\linewidth}{0.4mm}}\par\vspace{3mm}}
% A foundation-band header (Conclusions): centred, no index, twin hairline.
\newcommand{\baseband}[1]{%
  \begin{center}%
  {\sffamily\bfseries\fontsize{33}{37}\selectfont\color{char}#1}\\[2mm]
  {\color{teal}\rule{\linewidth}{0.4mm}}\end{center}\vspace{1.5mm}}

% Conference-grade header: large title on cream, then a full-width teal
% author band (white text) --- the standard academic-poster masthead
% (cf. Imperial/AOGS style). All author/venue metadata lives in the band.
\settitle{
  \begin{minipage}{0.99\linewidth}
  \centering
  {\color{char}\bfseries\fontsize{56}{63}\selectfont
   Discrete Layer Thermal Modeling of Lunar Landing Sites\\[6pt]
   with Topographic Shadowing}\\[7pt]
  % AOGS masthead rule: presenter's name UNDERLINED + BOLD; Abstract ID,
  % title, authors prominent; author e-mail + address provided.
  \colorbox{teal}{\parbox{0.985\linewidth}{\centering\vspace{3pt}
    {\color{white}\Large \underline{\textbf{Ram\'on III P.\ Gregorio}}$^{1}$%
     \enspace Richard Larsson$^{1}$ \enspace Yasuko Kasai$^{2}$}\\[6pt]
    {\color{white}\large $^{1}$Institute of Science Tokyo, Japan
     \enspace$\bullet$\enspace $^{2}$Institute of Science Tokyo, Transdisciplinary
     Science \& Engineering, School of Environment and Society, Japan}\\[4pt]
    {\color{white}\large \texttt{rp3gregorio@gmail.com} \enspace$\bullet$\enspace
     Abstract ID:\ [ID] \enspace$\bullet$\enspace Poster session:\ [window]}%
  \vspace{3pt}}}
  \end{minipage}
}

\begin{document}
\maketitle[titletotopverticalspace=2mm, titletoblockverticalspace=3mm]

\begin{columns}

% ════════════════════════════ LEFT COLUMN ═══════════════════════════════════
\column{0.5}

\block{}{
\horizon{1}{Physical model}
Lunar surface/subsurface temperatures need regolith thermal properties
\emph{and} topographically controlled sunlight; the only in-situ ground
truth is the Apollo 15/17 Heat Flow Experiment (Langseth et al.\ [2];
Nagihara et al.\ [3]). We add \textbf{two advances} to a certified 1-D
thermal-diffusion model (flux-anchored, $\sim$0.6 s/solve): \textbf{discrete
density layers} and \textbf{ray-traced shadowing} on the LOLA DEM [5].
\vspace{3pt}
\begin{tikzfigure}
  \includegraphics[width=0.94\linewidth]{../figures/poster/poster_schematic.pdf}
\end{tikzfigure}
\vspace{2pt}
Conductivity keeps the Hayne et al.\ [1] form
$K(T,z)=K_c(z)\,[1+\chi(T/350\,\mathrm{K})^3]$ with the flux-anchored
closure $\,d\langle T\rangle/dz=(Q_b-u_{\rm rect})/K\,$ below the
rectification zone. Two branches: \textbf{cp} --- density sets heat
capacity only, $K$ fixed at the certified
$K_d^{*}=\kdstarAfif{}/\kdstarAsev{}$ mW m$^{-1}$ K$^{-1}$; and
\textbf{coupled} --- density also sets $K_c$ through the model's own linear
$K_c(\rho)$, \emph{no new constants}. Both prefer a
\textbf{surface denser than Hayne's} (1300--1500 vs 1100 kg m$^{-3}$),
with the consolidated break \textbf{high} in the column.
\vspace{2pt}
\colorbox{teal!8}{\parbox{\linewidth}{\vspace{1pt}
\textbf{Topography enters twice:} the horizon cools the deep sensors
(\textbf{\shadowcoolAfif{}/\shadowcoolAsev{} K}), while the same terrain
radiates back ($+\svfdtirAfif{}$/$+\svfdtirAsev{}$ K IR;
$\mathrm{SVF}=\svfAfif{}/\svfAsev{}$) --- \textbf{two opposing first-order
channels}; cooling wins at A15, warming at A17.
\vspace{1pt}}}
}

\block{}{
\horizon{2}{The sites, their horizons, the lost sunlight}
\begin{tikzfigure}
  \includegraphics[width=0.62\linewidth]{../figures/poster/poster_sites.pdf}
\end{tikzfigure}
\vspace{2pt}
Apollo 15 sits under the Apennine front (horizon to \horAfif$^\circ$);
Apollo 17 inside Taurus--Littrow, its horizon peaking due north at the
North Massif (\horAsev$^\circ$). Shadowing deletes \lossAfif\% /
\lossAsev\% of the insolation energy --- the first and last slice of
every lunar day.
}

% ═══════════════════════════ RIGHT COLUMN ═══════════════════════════════════
\column{0.5}

\block{}{
\horizon{3}{Validation against Apollo HFE}
\begin{tikzfigure}
  \includegraphics[width=0.56\linewidth]{../figures/poster/poster_profiles.pdf}
\end{tikzfigure}
\vspace{2pt}
One equilibrium temperature per sensor from a quiet late-mission window
(1971--77 record; sensors above 80 cm excluded). The density-coupled winner
\emph{tilts} the deep gradient onto the sensors; the cp branch can only shift
it. All 650 configs converged; winners close the flux budget to
$\leq \closurewin\%$, agreeing to $\leq \driftmax$ mK at $n{=}96$.
\vspace{3pt}
\begin{tikzfigure}
  \includegraphics[width=0.46\linewidth]{../figures/poster/poster_transfer.pdf}
\end{tikzfigure}
\vspace{1pt}\par
{\small\color{char!75}\emph{Each site's structure moved to the other still
beats homogeneous --- physics, not curve-fitting.}}
}

\block{}{
\horizon{4}{Sensitivity: what actually matters}
\begin{tikzfigure}
  \includegraphics[width=0.50\linewidth]{../figures/poster/poster_sens.pdf}
\end{tikzfigure}
\vspace{2pt}
\textbf{Top:} heat capacity alone moves the fit $\leq 0.7$ K, while density
\textbf{through the conductivity it implies} spans ${\sim}4$ K.
\textbf{Bottom:} the mid$\vert$deep break sits \textbf{high} ($z_2 = 30$ cm) at
both sites. \textbf{The lever is conductivity, not heat capacity.}
}


\end{columns}

% ── full-width Conclusions band: read LAST, matching the abstract's arc ──────
% (Motivation -> method -> validation -> sensitivity -> cross-site ->
%  geology -> applications -> generalization). Two-column bullets keep it
%  compact across the full page.
\block{}{
\baseband{Robustness \& conclusions}
\colorbox{teal!8}{\parbox{\linewidth}{\vspace{2pt}
% left: the full pipeline (fills the band, recaps the method visually)
\begin{minipage}[t]{0.25\linewidth}
\centering
{\footnotesize\textbf{The full pipeline}\\[-1pt]
{\scriptsize\color{char!70} every structure runs the same chain}}\\[3pt]
\includegraphics[width=0.48\linewidth]{../figures/poster/poster_flowchart.pdf}
\end{minipage}\hfill
% middle: bootstrap robustness
\begin{minipage}[t]{0.29\linewidth}
\centering
\includegraphics[width=0.99\linewidth]{../figures/poster/poster_bootstrap.pdf}\\[2pt]
{\footnotesize\raggedright\textbf{10k-draw bootstrap:} both branches' best-RMSE
95\% CIs (coupled \bootcicKAfif{}/\bootcicKAsev{} K) sit far left of the
homogeneous baseline --- \emph{the improvement is robust}; the exact A15 triple survives
only \stabcKAfif\% of draws.\par}
\end{minipage}\hfill
% right: conclusions
\begin{minipage}[t]{0.42\linewidth}
\begin{itemize}
  \item \textbf{Topography is first-order} via two opposing channels:
        horizon blocking (\shadowcoolAfif{}/\shadowcoolAsev{} K) and
        $(1{-}\mathrm{SVF})$ re-radiation ($+\svfdtirAfif{}$/$+\svfdtirAsev{}$
        K; warms A17) --- same DEM horizons.
  \item A 3-layer column reproduces the record; \textbf{density matters
        through the conductivity it implies} --- coupled beats homogeneous
        ${\sim}2.6\times$ / $10\times$.
  \item Cross-site transfer \textbf{tracks the known geology} (physics, not
        fitting); extends to \textbf{landing-site assessment} and \textbf{any
        airless body}.
\end{itemize}
\end{minipage}
\vspace{2pt}}}

\vspace{3pt}
{\color{teal}\rule{\linewidth}{0.4mm}}\\[2pt]
{\footnotesize\textbf{References}\enspace
[1] Hayne, P.\ O., et al., J.\ Geophys.\ Res.\ Planets, vol.122, 2017,
pp.\ 2371--2400.\enspace
[2] Langseth, M.\ G., et al., Proc.\ Lunar Sci.\ Conf.\ 7, 1976,
pp.\ 3143--3171.\enspace
[3] Nagihara, S., et al., J.\ Geophys.\ Res.\ Planets, vol.123, 2018,
pp.\ 1125--1139.\enspace
[4] Saito, Y., et al., Adv.\ Space Res., vol.40, 2007, pp.\ 1601--1607.\enspace
[5] Smith, D.\ E., et al., Geophys.\ Res.\ Lett., vol.37, 2010, L18204.%
\hfill Code + archives: \texttt{github.com/rp3gregorio/Lunar-HFE}}
}

\end{document}
